\section{Introduction}

When two people stand side by side and we ask what they would look like as a single person, our
imagination instinctively tries to merge their features. We picture the eye shape of one person
combined with the jawline of the other, the nose of this person blended with the smile of that
person. This intuitive notion of ``averaging'' faces is precisely what computational face
averaging attempts to formalize. Yet what feels natural to the human brain becomes a surprisingly
rich algorithmic challenge when we ask a computer to perform this task.

This tutorial walks through every step of building a \emph{feature-based} face morphing system ---
one that produces convincing, natural-looking results by aligning facial geometry before blending
pixel values. Along the way, we will encounter concepts from computer vision, computational
geometry, linear algebra, and image processing. No single technique is particularly difficult, but
the way they fit together is elegant and instructive.


\subsection{What Is Face Averaging and Face Morphing?}

At its core, \emph{face averaging} answers a simple question: given two or more photographs of
different people's faces, can we produce a single image that represents their collective
appearance? The output should look like a plausible human face --- not a superimposed jumble, not
a half-and-half split, but a coherent image that could have been captured by a camera.

It helps to distinguish two closely related terms that are often used interchangeably but carry
subtly different meanings.

\begin{description}
    \item[Face Averaging] refers to the static problem of computing a single representative face
    from a set of input faces. The word ``average'' here is used in the everyday sense: we want
    something that embodies characteristics of all inputs. A common technique is to compute the
    mean of corresponding pixel values across aligned images, but as we will see shortly, a raw
    pixel average produces unsatisfactory results without careful geometric alignment first.

    \item[Face Morphing] refers to the dynamic, time-varying process of creating a smooth
    \emph{transition} from one face to another. A morph is a sequence of frames in which the
    first frame is face A, the last frame is face B, and every intermediate frame shows a
    progressively larger contribution from B and a smaller contribution from A. Face morphing
    is more general than face averaging because it requires producing not just one blended image
    but an entire continuous family of them, parameterized by a blending coefficient $\alpha \in
    [0, 1]$. When $\alpha = 0.5$, the morph frame is exactly an average of the two faces.
\end{description}

In this tutorial we focus primarily on the averaging problem ($\alpha = 0.5$ for two faces), but
everything we develop generalizes directly to morphing at arbitrary values of $\alpha$. Once we
know how to warp two faces into a common geometry and blend them at fifty percent, producing a
full morph sequence requires only a loop over different $\alpha$ values.

Formally, we can state the problem as follows.

\begin{definition}[Face Averaging]
Given $N$ input face images $I_1, I_2, \dots, I_N$, each containing a roughly frontal view of a
person's face, compute a single output image $I_{\text{avg}}$ that appears to be a plausible face
whose features are intermediate to those of all inputs. The output should satisfy two desirable
properties:
\begin{enumerate}
    \item \textbf{Geometric fidelity:} Facial features (eyes, nose, mouth, jawline) in
    $I_{\text{avg}}$ should have shapes and positions that are intermediate between the
    corresponding features in each $I_k$.
    \item \textbf{Photorealism:} $I_{\text{avg}}$ should look like a real photograph of a face,
    not an obvious composite or computer-generated image. There should be no visible seams,
    ghosting, or blurring artifacts caused by misalignment.
\end{enumerate}
\end{definition}

The distinction between these two properties is important. Geometric fidelity concerns
\emph{where} features are placed, while photorealism concerns \emph{how} they look once placed.
The difficulty of face averaging arises because these two concerns interact: even with
perfectly accurate pixel values, placing them at the wrong locations produces visible artifacts.
Conversely, even with perfect geometry, poor blending of color and texture yields unrealistic
results. Our algorithm must address both.


\subsection{Applications of Face Averaging and Morphing}

Face averaging and morphing are not merely academic curiosities. They appear in a wide variety of
real-world contexts, some beneficial and some concerning. Understanding the applications motivates
why we should care about producing high-quality results.


\subsubsection{Law Enforcement and Forensics}

One of the oldest and most socially valuable applications is in forensic imagery.
Investigating officers sometimes need to reconstruct what a missing person might look like after
years of absence, or to create a composite image from multiple partial descriptions.

\begin{itemize}
    \item \textbf{Age progression.} When a child has been missing for years, law enforcement
    agencies produce age-progressed images to show what the person might look like as an adult.
    While this typically involves more than simple averaging --- bone structure changes, skin
    texture evolves, hair recedes --- morphing techniques form a building block of the process.
    A sequence of morphed images at increasing $\alpha$ values can suggest how facial structure
    changes over time when combined with models of aging.

    \item \textbf{Composite sketches.} When multiple witnesses describe a suspect, their
    descriptions may differ. Averaging the facial features described by independent witnesses
    can produce a composite that captures the common elements while smoothing out inconsistencies
    in any single account.

    \item \textbf{Identifying unknown individuals.} In cases where only degraded or partial
    facial images are available, morphing can be used to enhance recognition by combining
    information from multiple sources.
\end{itemize}


\subsubsection{Entertainment and Visual Effects}

The entertainment industry has long been a driving force behind morphing technology. Some of the
most memorable visual effects in cinema history rely on face morphing.

\begin{itemize}
    \item \textbf{Film and television effects.} The 1992 music video for Michael Jackson's
    ``Black or White'' featured a famous morphing sequence where one person's face smoothly
    transforms into another's. Such effects, once requiring expensive proprietary hardware, can
    now be produced on a laptop using the algorithms described in this tutorial.

    \item \textbf{Celebrity morphs and deepfake precursors.} Interactive websites and applications
    that let users morph their own face with a celebrity's are perennially popular. While modern
    deepfake technology has largely moved to neural-network-based approaches, the traditional
    feature-based morphing techniques we study here remain the foundation on which much of this
    work is built.

    \item \textbf{Character design.} In animation and video game development, artists sometimes
    average multiple source faces to create a new character whose appearance feels familiar yet
    distinctive. Averages of many faces tend to produce unusually attractive results --- a
    phenomenon we will discuss below --- making this a useful tool for character design.
\end{itemize}


\subsubsection{Cognitive Science and Psychology}

Perhaps the most scientifically surprising application of face averaging comes from the study of
human face perception. In 1878, Francis Galton --- a pioneer of statistics --- had an idea that
turned out to be remarkably prescient. He proposed that by superimposing multiple photographic
portraits, one could produce a ``composite photograph'' that represents the typical appearance of
a group.

Galton's original method was crude: he simply exposed a single photographic plate multiple times,
once for each person's portrait, with each exposure reduced in intensity. The resulting composite
turned out to be more attractive than essentially any of the individual faces that went into it.
This observation led to decades of research.

Modern confirmations of Galton's finding include:

\begin{itemize}
    \item \textbf{The ``averageness'' hypothesis.} Studies have repeatedly shown that composite
    faces --- whether created by photographic averaging or by computer morphing --- are rated as
    more attractive than individual faces. One explanation is that averageness signals genetic
    diversity and health. Another is that average faces are easier for the visual system to
    process (a ``processing fluency'' account).

    \item \textbf{Prototype formation.} Human brains appear to form prototype representations of
    faces automatically. When we see many faces of a particular demographic group, our visual
    system constructs an implicit ``average'' of those faces, which then becomes the baseline
    against which individual faces are compared. Morphing experiments allow researchers to probe
    this process by presenting controlled intermediate faces and measuring human responses.

    \item \textbf{Own-race bias.} People tend to be better at recognizing faces of their own racial
    or ethnic group than faces of other groups. Morphing studies have helped clarify whether this
    bias arises from perceptual familiarity (the prototype for other-race faces is less well
    developed) or from social factors.
\end{itemize}


\subsubsection{Biometric Security and Morphing Attacks}

On the more concerning side of applications, face morphing plays an important role in understanding
and attacking biometric security systems.

A \emph{morphing attack} on a face recognition system works as follows. An attacker creates a
morphed image that combines their own face with the face of an authorized person (the ``target'').
Because the morph contains features of both individuals, it may be accepted by the enrollment
system as belonging to the authorized person. Once the morphed image is enrolled in the database,
the attacker can then authenticate using their \emph{own} face, because the system has learned to
associate features of the attacker's face with the target's identity.

This attack vector is particularly dangerous in border control applications. If a forged passport
contains a morphed photograph of two different people, both people could potentially use that
single passport. The European Union's border control system has identified this as a significant
vulnerability, and research into morph detection has become a distinct subfield.

Understanding morphing attacks requires understanding morphing algorithms at a deep level. This is
why we study the full pipeline in this tutorial: the same techniques that produce beautiful
entertainment morphs also produce the morphs that can fool biometric systems, and knowing
how they work is essential for building defenses.


\subsection{The Naive Approach: Why Simple Pixel Averaging Fails}

Before developing a sophisticated algorithm, it is instructive to understand why the obvious
approach does not work. Consider the following naive strategy:

\begin{enumerate}
    \item Take two face images $I_1$ and $I_2$.
    \item Align them by identifying the centers of the eyes in each image and applying a global
    similarity transform (rotation, translation, and uniform scaling) so that the eye centers match.
    \item For every pixel position $(x, y)$, compute the output pixel as the arithmetic mean:
    \[
        I_{\text{avg}}(x, y) = \frac{I_1(x, y) + I_2(x, y)}{2}.
    \]
\end{enumerate}

This approach is simple, requires no advanced mathematics, and can be implemented in a few lines
of code. So what goes wrong?

\subsubsection{The Ghosting Effect}

The fundamental problem is that two different faces, even when aligned by eye centers, do not agree
on the precise position of other facial features. Consider the nose. Person A may have a nose that
sits slightly lower on the face, or is wider, or has a different tip shape than Person B's nose.
After global eye-based alignment, the two noses will not occupy exactly the same set of pixels.

When we average pixel-by-pixel, the pixels belonging to A's nose and the pixels belonging to B's
nose are at slightly different locations. The result is that neither nose appears sharply in the
output --- instead, we see a blurry, ghost-like artifact where both noses are partially visible but
neither is crisp. The same issue occurs for every facial feature:

\begin{itemize}
    \item \textbf{Eyes} may have slightly different shapes and sizes; even when eye centers align,
    the corners of the eyes and the eyelids do not.
    \item \textbf{The nose} is often the worst offender, as nose position and shape vary
    considerably between individuals and the nose's shadow can appear in different places.
    \item \textbf{The mouth} varies in width, the corners may be at different vertical positions,
    and lip thickness differs significantly.
    \item \textbf{The jawline and face shape} often differ enough that the outline of the averaged
    face appears double-exposed.
    \item \textbf{Eyebrows} differ in thickness, arch position, and vertical placement.
    \item \textbf{Ears} may be visible in one image but not the other, or may appear at different
    positions relative to the face outline.
\end{itemize}

The cumulative effect is a blurred, ghostly image that looks like a double exposure of two faces
rather than a single coherent person. The blurring is not uniform: it is worst at the boundaries of
facial features, where the position disagreement is largest. This creates a particularly
unnatural appearance because the eyes and mouth --- which carry the most perceptually salient
information about a face --- are the regions where the artifacts are most visible.

\subsubsection{A Conceptual Example}

To make this concrete, imagine a one-dimensional version of the problem. Suppose we represent the
intensity profile along a horizontal line through an eye as a simple function. Person A's eye
extends from pixel 40 to pixel 60, while Person B's eye extends from pixel 43 to pixel 63 (it is
slightly shifted to the right). If we average these profiles pixel by pixel:

\begin{itemize}
    \item At pixels 40--42, only Person A contributes eye pixels; Person B contributes background
    (skin) pixels. The average is a weak, partial eye signal.
    \item At pixels 43--60, both contribute eye pixels, but at different sub-regions within the
    eye. The average is muddled.
    \item At pixels 61--63, the situation is reversed: Person B contributes eye pixels while Person
    A contributes background.
\end{itemize}

The result is a smeared eye that is wider than either original but lacks the crisp edges of either.
This is exactly what we observe in the two-dimensional case, only the misalignment occurs in both
horizontal and vertical directions with varying amounts at different locations on the face.

\subsubsection{Why Global Alignment Cannot Fix This}

The key insight is that the misalignment is \emph{local}, not global. A global similarity
transform (rotation, translation, uniform scale) can bring the eyes into alignment, but it cannot
simultaneously bring the nose, mouth, jawline, and all other features into alignment. These
features move independently of one another because they are governed by the underlying bone
structure, muscle, and soft tissue, which differ between individuals.

What we need is a \emph{local} deformation that can independently adjust different regions of
the face. This is the essence of the feature-based approach: instead of asking for a single
transformation of the entire image, we divide the face into regions and compute a separate
transformation for each region, ensuring that corresponding features are brought into precise
alignment before any pixel averaging occurs.


\subsection{The Feature-Based Solution: A High-Level View}

The feature-based approach to face averaging solves the ghosting problem by inverting the order of
operations. Rather than averaging pixel values first and hoping the geometry works out, we ensure
the geometry is correct first, then average.

The idea proceeds in three conceptual steps.

\paragraph{Step 1: Establish correspondences.}
We identify a set of \emph{corresponding points} in each input face. These are landmark points
that have the same semantic meaning in every face: the tip of the nose, the corner of the left
eye, the center of the lower lip, and so on. If we identify 68 such points in each face, then
point $i$ in Face A corresponds to point $i$ in Face B for all $i = 1, \dots, 68$.

\paragraph{Step 2: Define a target geometry.}
Given the landmarks in each face, we compute a \emph{mean shape} by averaging the landmark
positions. For 68 points, this gives us 68 mean coordinates. The mean shape represents the
``average'' placement of features and will serve as the geometric canvas onto which we project
each face.

\paragraph{Step 3: Warp and blend.}
We deform (warp) each input face so that its landmarks land exactly on the mean shape's
landmarks. After warping, corresponding features in all faces occupy the same pixel locations.
Only then do we average the pixel values. Because the geometry is aligned, the averaging produces
a sharp, natural-looking result rather than a ghostly blur.

This approach is called \emph{feature-based} because the deformation is driven by the positions
of specific facial features (the landmarks), and the intermediate geometry (the triangulation and
warping) is constructed to respect those feature locations.

\subsubsection{Piecewise Affine Warping}

The warping step requires care. We cannot simply warp the entire face with a single transformation
because different parts of the face need to move by different amounts. The nose might need to
shift downward while the mouth shifts upward. Instead, we divide the face into a mesh of triangles
using the landmark points as vertices. Each triangle undergoes its own affine transformation.

An \emph{affine transformation} preserves straight lines and parallelism. It can translate,
rotate, scale, and shear, but it cannot bend or curve. By restricting ourselves to affine
transforms within each triangle, we ensure that the local structure of the image within the
triangle is preserved while the triangle as a whole moves to its target position.

This approach is called \emph{piecewise affine warping} because the overall transformation is
a patchwork of affine transforms, one per triangle, joined together at the triangle boundaries.
The result is a continuous deformation of the image that precisely maps each landmark to its
target position.

The use of triangles is not arbitrary. Triangles have several desirable properties:

\begin{itemize}
    \item Three points uniquely define a triangle, and three point correspondences uniquely define
    an affine transformation. This makes the mathematics straightforward.
    \item Triangles tile any planar point set (a property not shared by, say, quadrilaterals).
    Using the Delaunay triangulation of the landmark points, we can decompose the entire face
    region into non-overlapping triangles with good geometric properties.
    \item The Delaunay triangulation maximizes the minimum angle of all triangles, avoiding
    long, skinny triangles that would produce poor interpolation results.
\end{itemize}


\subsection{Algorithm Pipeline}

The complete face averaging pipeline, from raw input images to a blended average face, consists of
the following five major steps. Each step is the subject of a subsequent section in this tutorial,
but we present the full pipeline here so you can see the big picture before diving into details.

\begin{enumerate}
    \item \textbf{Detect facial landmarks.} Given an input face image, we run a landmark detector
    to find the coordinates of key facial features. A typical detector returns 68 points: the jaw
    line (17 points), the left and right eyebrows (5 points each), the nose bridge and nose tip (9
    points), the left and right eyes (6 points each), and the mouth outline (20 points). Popular
    detectors include dlib's shape predictor and MediaPipe's FaceMesh. The choice of detector
    affects the number and precision of landmarks but not the overall algorithm structure.

    \item \textbf{Compute mean landmark positions.} With landmarks identified in each input face,
    we compute the component-wise arithmetic mean of corresponding landmark positions. If Face A
    has landmark $i$ at position $\mathbf{p}^{(A)}_i = (x^{(A)}_i, y^{(A)}_i)$ and Face B
    has landmark $i$ at $\mathbf{p}^{(B)}_i = (x^{(B)}_i, y^{(B)}_i)$, then the mean position is:
    \[
        \mathbf{p}^{(\text{mean})}_i = \frac{\mathbf{p}^{(A)}_i + \mathbf{p}^{(B)}_i}{2}.
    \]
    For $N$ faces, the mean is simply the sum over all $N$ divided by $N$. The resulting set of
    mean landmark positions defines the mean shape.

    \item \textbf{Triangulate the mean shape.} Using the mean landmark positions as vertices, we
    compute a Delaunay triangulation. This partitions the convex hull of the landmarks into a set
    of triangles such that no landmark point lies inside the circumcircle of any triangle. The
    Delaunay triangulation is uniquely determined by the point set and has the important property
    of maximizing the minimum triangle angle, which avoids degenerate (long, thin) triangles that
    would cause numerical instability during the warping step.

    \item \textbf{Warp each face to the mean shape.} For each input face and each triangle in the
    Delaunay triangulation of the mean shape, we perform the following sub-steps:
    \begin{enumerate}
        \item[a.] Find the corresponding triangle in the input face (using the same landmark
        indices).
        \item[b.] Compute the affine transformation that maps the input face's triangle to the
        mean shape's triangle.
        \item[c.] Apply this transformation (in reverse, via inverse warping) to copy pixel values
        from the input face into the target triangle in the mean shape.
    \end{enumerate}
    This process is repeated for every triangle, yielding a complete warped image of each face in
    the mean shape geometry. The warping uses inverse mapping and interpolation to avoid gaps and
    ensure every output pixel gets a value.

    \item \textbf{Blend the warped faces.} Once all faces are warped to the common mean geometry,
    their pixel values can be directly averaged. For a two-face morph at blending coefficient
    $\alpha$, the output image is:
    \[
        I_{\text{out}} = (1 - \alpha) \cdot I^{(A)}_{\text{warped}} + \alpha \cdot
        I^{(B)}_{\text{warped}}.
    \]
    For face averaging with equal weights, $\alpha = 0.5$. For a full morph sequence, we compute
    this expression for multiple values of $\alpha$ from 0 to 1, producing a smooth transitions
    from Face A to Face B.
\end{enumerate}

The following figure illustrates the complete pipeline as a flow diagram.

\begin{figure}[htbp]
    \centering
    \begin{tikzpicture}[
        >=Stealth,
        node distance=1.6cm,
        box/.style={
            rectangle,
            rounded corners=3pt,
            draw=blue!60!black,
            fill=blue!5,
            minimum width=4.2cm,
            minimum height=1.2cm,
            align=center,
            font=\small
        },
        inputbox/.style={
            rectangle,
            rounded corners=3pt,
            draw=orange!80!black,
            fill=orange!10,
            minimum width=3.5cm,
            minimum height=1.0cm,
            align=center,
            font=\small
        },
        outputbox/.style={
            rectangle,
            rounded corners=3pt,
            draw=green!60!black,
            fill=green!5,
            minimum width=3.5cm,
            minimum height=1.0cm,
            align=center,
            font=\small
        },
        stepnum/.style={
            circle,
            draw=blue!60!black,
            fill=blue!60!black,
            text=white,
            minimum size=0.6cm,
            font=\footnotesize\bfseries,
            inner sep=0pt
        },
        arrow/.style={
            ->,
            thick,
            blue!60!black
        }
    ]
    % Title
    \node[font=\large\bfseries, text=blue!60!black] (title)
        {Face Averaging Pipeline};

    % Input faces (parallel branches)
    \node[inputbox, below=1.5cm of title, xshift=-2.5cm] (faceA)
        {\textbf{Face Image A}\\$I_A$};
    \node[inputbox, below=1.5cm of title, xshift=2.5cm] (faceB)
        {\textbf{Face Image B}\\$I_B$};

    % Step 1: Landmark Detection
    \node[box, below=1.8cm of title, xshift=-2.5cm] (step1a)
        {Detect 68 Landmarks};
    \node[box, below=1.8cm of title, xshift=2.5cm] (step1b)
        {Detect 68 Landmarks};
    \node[stepnum, at={($(faceA)!0.5!(step1a)$)}] (s1label) {1};

    % Landmark outputs
    \node[box, below=of step1a, font=\footnotesize] (landA)
        {Landmarks $\{\mathbf{p}^{(A)}_i\}$};
    \node[box, below=of step1b, font=\footnotesize] (landB)
        {Landmarks $\{\mathbf{p}^{(B)}_i\}$};

    % Step 2: Mean Shape
    \node[box, below=1.5cm of landA, xshift=2.5cm] (step2)
        {Compute Mean Shape\\$\mathbf{p}^{(\text{mean})}_i = \frac{1}{N}\sum_k
        \mathbf{p}^{(k)}_i$};

    % Step 3: Triangulation
    \node[box, below=of step2] (step3)
        {Delaunay Triangulation\\of Mean Shape};

    % Step 4: Warping (two parallel paths)
    \node[box, below=1.5cm of step3, xshift=-2.5cm] (step4a)
        {Warp Face A\\to Mean Shape\\(Piecewise Affine)};
    \node[box, below=1.5cm of step3, xshift=2.5cm] (step4b)
        {Warp Face B\\to Mean Shape\\(Piecewise Affine)};
    \node[stepnum, left=3mm of step4a] () {4};

    \node[box, below=0.8cm of step4a, font=\footnotesize] (warpA)
        {$I_A^{\text{warped}}$};
    \node[box, below=0.8cm of step4b, font=\footnotesize] (warpB)
        {$I_B^{\text{warped}}$};

    % Step 5: Blending
    \node[box, below=1.0cm of warpA, xshift=2.5cm] (step5)
        {Alpha Blending\\$I_{\text{out}} = (1-\alpha) I_A^w + \alpha I_B^w$};

    % Output
    \node[outputbox, below=of step5] (output)
        {\textbf{Average Face}\\$I_{\text{avg}}$};

    % Arrows: Face -> Landmark Detection
    \draw[arrow] (faceA) -- (step1a);
    \draw[arrow] (faceB) -- (step1b);

    % Arrows: Landmark Detection -> Landmark outputs
    \draw[arrow] (step1a) -- (landA);
    \draw[arrow] (step1b) -- (landB);

    % Arrows: Landmarks -> Mean Shape
    \draw[arrow] (landA) -- (landA |- step2.north);
    \draw[arrow] (landB) -- (landB |- step2.north);

    % Arrow: Mean Shape -> Triangulation
    \draw[arrow] (step2) -- (step3);

    % Arrows: Triangulation -> Warping
    \draw[arrow] (step3.south) -- +(0, -0.5cm) -| (step4a.north);
    \draw[arrow] (step3.south) -- +(0, -0.5cm) -| (step4b.north);

    % Arrow: Face A -> Warp A (implicit: use same triangulation)
    \draw[arrow, dashed, orange!80!black] (faceA.west) -- ++(-1.0cm, 0) |- (step4a.west)
        node[midway, left, font=\tiny, text=orange!80!black] {image data};
    \draw[arrow, dashed, orange!80!black] (faceB.east) -- ++(1.0cm, 0) |- (step4b.east)
        node[midway, right, font=\tiny, text=orange!80!black] {image data};

    % Arrows: Warped -> Blending
    \draw[arrow] (warpA) -- (warpA |- step5.north);
    \draw[arrow] (warpB) -- (warpB |- step5.north);

    % Arrow: Blending -> Output
    \draw[arrow] (step5) -- (output);

    \end{tikzpicture}
    \caption{Complete face averaging pipeline. Two input face images (orange, top) are processed
    in parallel through landmark detection. The extracted landmark sets are combined to compute a
    mean shape (Step~2), which is then triangulated using Delaunay triangulation (Step~3).
    The triangulation guides piecewise affine warping of each face into the mean geometry
    (Step~4). Finally, the warped images are blended via alpha compositing (Step~5) to produce
    the output average face (green, bottom). Dashed orange arrows indicate that the original
    image pixel data flows directly into the warping stage, guided by the geometric structure
    computed in Steps 2 and 3.}
    \label{fig:pipeline}
\end{figure}

\subsection{Why This Approach Produces Natural Results}

To close this introduction, it is worth understanding at an intuitive level why the feature-based
approach succeeds where the naive approach fails.

The core reason is \emph{structural alignment}. When a human observer looks at a face, the brain
does not process it pixel by pixel. Instead, it identifies a structure: there are two eyes, a nose,
a mouth, arranged in a particular configuration. The recognition and interpretation of the face
depends on getting this structure right.

The naive pixel-averaging approach operates at the level of the retina (the pixel grid), ignoring
structure entirely. It assumes that pixel $(x, y)$ in one face corresponds to pixel $(x, y)$ in
the other face, which is simply not true for any features beyond the alignment anchors.

The feature-based approach, by contrast, operates at the level of the structure. By identifying
corresponding landmarks and warping to a common geometry, we ensure that the left eye of Face A
and the left eye of Face B are mapped to the same location in the output before their pixel
values are averaged. The structure is aligned first, so the pixel averaging makes sense.

The result is a face that looks like a real person because the internal structure is internally
consistent. All features are in their proper relative positions, just as they would be in a
photograph. The eyes are sharp because both input eye images have been warped to the same
geometry before averaging, rather than being averaged at slightly offset positions. The same is
true for the nose, mouth, and all other features.

This structural perspective also explains why 68 landmarks (or similar) are sufficient to produce
convincing results even though a typical face image contains millions of pixels. The landmarks
capture the essential structural degrees of freedom of the face --- the positions of eyes, nose,
mouth, jaw, and eyebrows --- and the piecewise affine warping fills in the continuous deformation
between these control points. The pixel values carry the texture and color information, but the
structure is entirely determined by the landmarks.

In the following sections, we will examine each of these five pipeline steps in detail, working
through the mathematics, the implementation, and the design choices that go into building a
complete face morphing system.


\endinput
